\section{Conclusion and Future Work}\label{sec:conclusion}

This project implemented a TF-IDF Vector Space retrieval system for the
Cranfield dataset and evaluated it using standard ranked retrieval metrics. The
final system uses Punkt sentence segmentation, Penn Treebank tokenization,
Porter stemming, stemmed stopword removal, smoothed TF-IDF, cosine similarity,
an inverted TF-IDF index, and Rocchio-style pseudo-relevance feedback.

The improved system outperformed the no-feedback baseline at $k=10$ for
Precision, Recall, F0.5, MAP, and nDCG. The strongest summary result is the
increase in MAP@10 from 0.3247 to 0.3639, supported by a paired sign test over
per-query AP@10 values. Therefore, for the Cranfield ranked retrieval task, the
implemented feedback-enhanced TF-IDF system is better than the basic TF-IDF
cosine baseline with respect to MAP@10, under the assumption that the initial
top-ranked document is usually useful for expansion.

Future work should address the remaining weaknesses. First, spelling correction
or character n-gram matching could reduce zero-result failures caused by
out-of-vocabulary and misspelled query terms. Second, feedback parameters
$(m,\alpha,\beta)$ should be tuned on a validation split instead of fixed
manually. Third, semantic methods such as LSA could be compared against the
current feedback model to test whether latent concepts improve retrieval without
causing excessive query drift.
